\chapter{The solution method}\label{chap:solution}

\manref{cge_solve} solves the system of equations of
chapter~\ref{chap:equations} under the closure of chapter~\ref{chap:closures}.
The variables of all blocks are stacked in one vector. Variables that do not
exist for an element (for example the imports of a commodity that is not
imported) stay in the vector at zero and are left out of the system, and so
are the variables that the closure fixes. The free variables $x$ and the
equations $r(x) = 0$ form a square system.

\section{Newton's method}\label{sec:newton}

The solver works on scaled variables and scaled residuals,
\begin{equation*}
z_j = x_j / s_j, \qquad F_i(z) = r_i(x) / e_i,
\end{equation*}
where $s_j$ and $e_i$ are the base magnitudes of the variable and of the
equation (one where the base value is zero), so that the tolerance applies
alike to every block. Starting from the base, each Newton step solves
\begin{equation*}
J(z)\,\Delta z = -F(z)
\end{equation*}
by the LU factorisation of the Jacobian $J$, and the new point is
$z + \lambda\,\Delta z$, with $\lambda$ the largest of $1, \tfrac12, \tfrac14,
\dots$ (down to $10^{-8}$) such that
\begin{equation*}
\lVert F(z + \lambda\,\Delta z)\rVert \le (1 - 10^{-4}\lambda)\,\lVert F(z)\rVert,
\end{equation*}
a backtracking line search on the Euclidean norm of the residuals. Points at
which a price or quantity is not positive are rejected, and the step is
shortened. The solve has converged when $\max_i |F_i| < \tau$, the tolerance
(\texttt{1e-10} by default), over every equation of the model, including the
definitions left out of the reduced system (section~\ref{sec:reduced}).

The LU factors of a Jacobian are reused for further steps while each full
step at least halves $\lVert F\rVert$, a chord method; otherwise, or after a
step shortened by the line search, a new Jacobian is computed. Near a
solution Newton's method with a new Jacobian at each step converges
quadratically, and the chord method linearly; each chord step costs one
evaluation of the equations and two triangular solves, against a Jacobian and
a factorisation for a Newton step, so for large models it is much faster.

\section{The reduced system}\label{sec:reduced}

Some equations are explicit definitions: a variable equals an expression in
other variables, and the variable appears in none of the equations whose rows
have exact derivatives (section~\ref{sec:jacobian}). Such definitions are left
out of the Newton system, and their variables are computed from the other
variables before each evaluation of the residuals. They are
\eqref{eq:P1}, \eqref{eq:P2} and \eqref{eq:P6} (import and export prices and
the price of intermediate input), \eqref{eq:Q5} (intermediate demand),
\eqref{eq:Q6} (output by commodity), \eqref{eq:I5} (household demand),
\eqref{eq:T8} (demand for margin services), \eqref{eq:I13} (the purchases of
the health insurance fund) and \eqref{eq:H2} (factor efficiency). A definition
is left out only when its equation elements are exactly the existing elements
of its variable. Variables computed from their definitions cannot be fixed by
the closure (\secref{sec:swaps}).

The left-out equations still hold exactly at every point. Intermediate and
household demand grow with the square of the number of sectors, so most of a
large model's equations are left out: for a SAM with 250 activities and 20
households, 4,254 of 40,190 equations remain. \cmd{cge_solve, full} keeps
every equation; the test suite checks that both systems give the same
solution. The pattern of the reduced system is found by substituting each
left-out variable by the variables of its definition.

\section{The Jacobian}\label{sec:jacobian}

The equations that involve many variables, the price indexes \eqref{eq:P8}
and \eqref{eq:P9}, factor income \eqref{eq:I1}, government revenue and
spending \eqref{eq:I8} and \eqref{eq:I9}, factor markets \eqref{eq:S1}, the
balance of payments \eqref{eq:S3}, the savings--investment balance
\eqref{eq:S5}, commodity markets \eqref{eq:S2} and activity prices
\eqref{eq:P7}, have exact derivatives. In the reduced system these are total
derivatives through the left-out variables. The other rows are sparse. Their
pattern follows from the equations and is built at calibration, and their
derivatives are finite differences taken over groups of columns that share no
row \citep{curtis1974}: the columns are coloured greedily, the most connected
first, so that one evaluation of the equations gives the derivatives of every
column of a group. For a group $G$, with $e_G$ the vector with ones in the
columns of $G$,
\begin{equation*}
J_{ij} \approx \frac{F_i(z + h \odot e_G) - F_i(z)}{h_j}, \qquad j \in G,
\qquad h_j = 10^{-7} \max(1, |z_j|),
\end{equation*}
one-sided differences that use the residuals at $z$, which Newton's method has
already computed. Their error, measured against column-by-column central
differences, is between $3 \times 10^{-7}$ and $10^{-6}$ in the test models,
too small to change the Newton steps in those models; the test suite checks
that Newton's method reaches the same solution with both kinds of
differences. When Newton's method fails with them it
is tried again from the start with central differences,
\begin{equation*}
J_{ij} \approx \frac{F_i(z + h \odot e_G) - F_i(z - h \odot e_G)}{2 h_j}, \qquad h_j = 10^{-6} \max(1, |z_j|),
\end{equation*}
whose error is about $10^{-10}$ but which need two evaluations of the
equations for each group. The diagnostics (section~\ref{sec:diagnostics}) use
central differences.

On a 240-account SAM with 4,439 equations, computing the Jacobian column by
column took about 250 seconds per Newton iteration in Stata 13/IC, against a
dense LU solve of under 3 seconds; the grouped Jacobian takes a few seconds.
Most of the time of a large solve is spent evaluating the equations for the
Jacobian (\secref{sec:performance}).

\section{Fallbacks and continuation}\label{sec:continuation}

If Newton's method fails with both kinds of differences, Powell's hybrid
method is tried (Mata's \cmd{solvenl()}, technique \cmd{broydenpowell}),
followed by Newton steps to polish the result.

If that fails too, the shock is applied by continuation: the model is solved
for $k = 1, \dots, n$ with parameters and fixed values
\begin{equation*}
\theta_k = \theta^0 + \frac{k}{n}\,(\theta^1 - \theta^0),
\end{equation*}
between their base values $\theta^0$ and their shocked values $\theta^1$,
each step starting from the solution of the step before; by default in
$n = 4$ steps and then in $n = 16$. A shock that defeats continuation is often
one that the economy cannot absorb.

\section{Diagnostics}\label{sec:diagnostics}

After a failure, or on request, \manref{cge_solve} displays:
\begin{itemize}
\item the equations with the largest scaled residuals at the last point, by
block and element;
\item any equation with no free variable, and any free variable in no
equation, which make the Jacobian singular;
\item the smallest pivot of the LU factorisation of the Jacobian relative to
the largest, and the variable it belongs to: partial pivoting permutes rows
only, so column $j$ of $U$ belongs to variable $j$.
\end{itemize}

\section{Constraints}\label{sec:solution-constraints}

With constraints (chapter~\ref{chap:constraints}) the model is solved in two
stages: first with smoothed complementarity conditions, which find whether each
constraint binds, and then with the exact conditions in the regimes found,
changing the regime of any constraint whose other condition fails, for at most
ten rounds (\secref{sec:cons-solution}).

\section{Other implementations}

The Python reference implementation uses the same Newton method with a
backtracking line search, and its own pure active-set method for constraints.
The GAMS program written by \manref{cge_gams} is solved by the solvers that
GAMS provides, PATH by default (chapter~\ref{chap:gams}). Their solutions agree
with those of \cmd{cge_solve} to the precision of the solutions
(chapter~\ref{chap:validation}).
